\documentclass[aip,jcp,reprint,noshowkeys,superscriptaddress]{revtex4-1}
\usepackage{graphicx,dcolumn,bm,xcolor,microtype,multirow,amscd,amsmath,amssymb,amsfonts,physics,longtable,wrapfig}
\usepackage[version=4]{mhchem}

\usepackage[utf8]{inputenc}
\usepackage[T1]{fontenc}
\usepackage{txfonts}

\usepackage[
    colorlinks,
    linkcolor={red!50!black},
    citecolor={red!70!black},
    urlcolor={red!80!black}
	]{hyperref}
\urlstyle{same}

\newcommand{\ie}{\textit{i.e.}}
\newcommand{\eg}{\textit{e.g.}}
\newcommand{\alert}[1]{\textcolor{red}{#1}}
\usepackage[normalem]{ulem}
\newcommand{\titou}[1]{\textcolor{red}{#1}}
\newcommand{\trashPFL}[1]{\textcolor{\red}{\sout{#1}}}
\newcommand{\PFL}[1]{\titou{(\underline{\bf PFL}: #1)}}

\newcommand{\cre}[1]{\Hat{a}_{#1}^{\dag}}
\newcommand{\ani}[1]{\Hat{a}_{#1}}
\newcommand{\ca}[2]{\Hat{a}^{#1}_{#2}}

\newcommand{\mc}{\multicolumn}
\newcommand{\fnm}{\footnotemark}
\newcommand{\fnt}{\footnotetext}
\newcommand{\tabc}[1]{\multicolumn{1}{c}{#1}}
\newcommand{\QP}{\textsc{quantum package}}
\newcommand{\T}[1]{#1^{\intercal}}

% coordinates
\newcommand{\br}{\boldsymbol{r}}
\newcommand{\bx}{\boldsymbol{x}}
\newcommand{\bI}{\boldsymbol{1}}
\newcommand{\cK}{\mathcal{K}}

% orbital energies
\newcommand{\eps}{\epsilon}
\newcommand{\irr}{\text{irr}}

% shortcuts for greek letters
\newcommand{\si}{\sigma}
\newcommand{\la}{\lambda}
\newcommand{\ii}{\mathrm{i}}
\newcommand{\up}{\uparrow}
\newcommand{\dw}{\downarrow}
\newcommand{\upup}{\up\up}
\newcommand{\updw}{\up\dw}
\newcommand{\dwup}{\dw\up}
\newcommand{\dwdw}{\dw\downarrow}

\newcommand{\h}{\text{1h}}
\newcommand{\p}{\text{1p}}
\newcommand{\hp}{\text{1h+1p}}
\newcommand{\hhp}{\text{2h1p}}
\newcommand{\pph}{\text{2p1h}}
\newcommand{\ph}{\text{ph}}
\newcommand{\eh}{\text{eh}}
\newcommand{\he}{\text{he}}
\newcommand{\ee}{\text{ee}}
\newcommand{\teh}{\overline{\text{eh}}}
\newcommand{\pp}{\text{pp}}
\newcommand{\hh}{\text{hh}}
\newcommand{\xc}{\text{xc}}
\newcommand{\Hx}{\text{Hx}}

% addresses
\newcommand{\LCPQ}{Laboratoire de Chimie et Physique Quantiques (UMR 5626), Universit\'e de Toulouse, CNRS, UPS, France}

\begin{document}

\title{Algebraic Reduction and Linearization of the Three-Frequency Bethe--Salpeter Equation through Third Order}

\author{Pierre-Fran\c{c}ois \surname{Loos}}
	\email{loos@irsamc.ups-tlse.fr}
	\affiliation{\LCPQ}

\begin{abstract}
The exact electron-hole Bethe--Salpeter equation is an integral equation in two fermionic frequencies in addition to the bosonic transfer frequency.
These notes examine when this genuine three-frequency problem can be reduced to a one-frequency algebraic eigenvalue problem without changing its perturbative excitation poles.
Using the exact factorization of the noninteracting electron-hole pair propagator, the transition amplitude is separated into the usual $Q^\pm$ frequency shape and a complementary frequency-dependent component.
The feedback of this complementary component is analyzed explicitly at second order in the interaction.
After centering the one-sided external vertices, same-sector $\hhp$--$\hhp$ and $\pph$--$\pph$ covariances vanish identically, while the mixed $\hhp$--$\pph$ covariance contains a factor that cancels the potentially resonant one-particle--one-hole propagator.
Under the usual non-intruder assumptions, the first contribution of the discarded frequency-shape sector to the excitation pole therefore occurs at fourth order.
Consequently, the two-sided projected Bethe--Salpeter equation is perturbatively equivalent to the full three-frequency pole equation through third order.
The projected kernel is then organized into static terms, neutral-double resolvents, second-order corrections to the external couplings, and a first-order interaction inside the neutral-double space.
This rational one-frequency kernel can be linearized exactly through third order as a frequency-independent generalized eigenvalue problem in an enlarged auxiliary space.
A spectator-labelled construction is introduced for the third-order external residues, together with a compact recipe for constructing and checking the resulting algebraic problem.
\end{abstract}

\maketitle

\onecolumngrid
\tableofcontents
\twocolumngrid

%%%%%%%%%%%%%%%%%%%%%%%%%
\section{Scope and notation}
\label{sec:scope}
%%%%%%%%%%%%%%%%%%%%%%%%%

The objective is to connect two descriptions of neutral excitations.
The first is the exact three-frequency electron-hole Bethe--Salpeter equation (BSE), in which the irreducible vertex depends on two fermionic frequencies and one bosonic transfer.
The second is an algebraic Casida-like eigenvalue problem containing ordinary frequency-independent matrices and auxiliary resolvents.
The construction is developed perturbatively through third order in the two-electron interaction $v$ using a Hartree--Fock (HF) reference and a M\o ller--Plesset partition.

The direct $\eh$ convention is used throughout.
The four external fermionic line energies are
\begin{equation}
	\nu
	\qquad
	\nu+\omega
	\qquad
	\nu'
	\qquad
	\nu'+\omega
	\label{eq:external_frequencies}
\end{equation}
The transverse $\eh$ and particle-particle transfers are
\begin{align}
	\Omega_{\teh}
	&=
	\nu'-\nu
	\\
	\Omega_\pp
	&=
	\omega+\nu+\nu'
	\label{eq:channel_transfers}
\end{align}
The parquet decomposition relevant to the direct $\eh$ irreducible vertex is
\begin{equation}
	\Gamma^\eh
	=
	\Lambda
	+
	\Phi^{\teh}
	+
	\Phi^\pp
	\label{eq:Gamma_parquet}
\end{equation}
For a two-body interaction,
\begin{equation}
	\Lambda
	=
	v
	+
	\mathcal O(v^4)
	\label{eq:Lambda_order}
\end{equation}
so that through third order all nontrivial corrections to $\Gamma^\eh$ are generated by the transverse $\eh$ and $\pp$ reducible vertices.

It is convenient to absorb the conventional factor of $\ii$ into the BSE kernel
\begin{equation}
	\mathcal K
	=
	\ii\Gamma^\eh
	\label{eq:K_def}
\end{equation}
and to use the composite pair index
\begin{equation}
	\alpha
	=
	(ia,\sigma)
	\qquad
	\sigma=\pm
	\label{eq:alpha_def}
\end{equation}
The labels $+$ and $-$ denote forward and backward $\eh$ propagation, while $ia$ remains the orbital-pair label in both sectors.

%%%%%%%%%%%%%%%%%%%%%%%%%
\section[Exact three-frequency BSE]{Exact three-frequency BSE and one-sided factors}
\label{sec:exact_bse}
%%%%%%%%%%%%%%%%%%%%%%%%%

The exact correlation function obeys
\begin{equation}
	L
	=
	L_0
	+
	L_0\Gamma^\eh L
	\label{eq:BSE_symbolic}
\end{equation}
which in frequency space contains two internal fermionic-frequency integrations
\begin{equation}
	\begin{split}
	L(\nu,\nu',\omega)
	={}&
	L_0(\nu,\nu',\omega)
	\\
	&+
	\frac{1}{(2\pi)^2}
	\int d\bar\nu\,d\bar\nu'
	L_0(\nu,\bar\nu,\omega)
	\Gamma^\eh(\bar\nu,\bar\nu',\omega)
	L(\bar\nu',\nu',\omega)
	\end{split}
	\label{eq:BSE_three_frequency}
\end{equation}
For occupied and virtual HF orbitals,
\begin{align}
	G_{0,i}(\nu)
	&=
	\frac{1}{\nu-\eps_i-\ii\eta}
	&
	G_{0,a}(\nu)
	&=
	\frac{1}{\nu-\eps_a+\ii\eta}
	\label{eq:G0}
\end{align}
The forward and backward one-frequency pair propagators are
\begin{align}
	(L_0^+)_{ia}(\omega)
	&=
	\frac{\ii}
	{\omega-(\eps_a-\eps_i)+2\ii\eta}
	\\
	(L_0^-)_{ia}(\omega)
	&=
	-\frac{\ii}
	{\omega+(\eps_a-\eps_i)-2\ii\eta}
	\label{eq:L0_pm}
\end{align}
Define
\begin{align}
	Q^+_{ia}(\nu,\omega)
	&=
	-\ii
	\left[
	G_{0,i}(\nu)
	-
	G_{0,a}(\nu+\omega)
	\right]
	\\
	Q^-_{ia}(\nu,\omega)
	&=
	+\ii
	\left[
	G_{0,a}(\nu)
	-
	G_{0,i}(\nu+\omega)
	\right]
	\label{eq:Q_pm}
\end{align}
Then the three-frequency noninteracting pair propagator factorizes exactly
\begin{equation}
	L_{0,ia}^{\sigma}(\nu,\nu',\omega)
	=
	2\pi\delta(\nu-\nu')
	Q^\sigma_{ia}(\nu,\omega)
	(L_0^\sigma)_{ia}(\omega)
	\label{eq:L0_factorization}
\end{equation}
with
\begin{equation}
	\frac{1}{2\pi}
	\int d\nu\,
	Q^\sigma_{ia}(\nu,\omega)
	=
	1
	\label{eq:Q_normalization}
\end{equation}
and
\begin{equation}
	Q^-_{ia}(\nu,\omega)
	=
	Q^+_{ia}(\nu+\omega,-\omega)
	\label{eq:Q_reversal}
\end{equation}
In the remainder, integer multiples of $\eta$ are written as $0^+$ when only the causal prescription matters.

Near a neutral pole $\Omega_n$, write
\begin{equation}
	L_{\alpha\beta}(\nu,\nu',\omega)
	\underset{\omega\to\Omega_n}{\sim}
	\frac{
	\mathcal X_\alpha(\nu)
	\overline{\mathcal X}_\beta(\nu')
	}{
	\omega-\Omega_n
	}
	\label{eq:pole_form}
\end{equation}
Using Eq.~\eqref{eq:L0_factorization}, the homogeneous BSE for the transition amplitude becomes
\begin{equation}
	\mathcal X_\alpha(\nu)
	=
	Q_\alpha(\nu,\Omega)
	(L_0)_\alpha(\Omega)
	\sum_\beta
	\int\frac{d\nu'}{2\pi}
	\mathcal K_{\alpha\beta}(\nu,\nu',\Omega)
	\mathcal X_\beta(\nu')
	\label{eq:pole_equation}
\end{equation}
No frequency projection has been made in Eq.~\eqref{eq:pole_equation}.

%%%%%%%%%%%%%%%%%%%%%%%%%
\section[Kernel projections]{One-sided and two-sided kernel projections}
\label{sec:projections}
%%%%%%%%%%%%%%%%%%%%%%%%%

The exact one-sided projections of the irreducible kernel are
\begin{equation}
	U_{\alpha\beta}(\nu,\omega)
	=
	\frac{1}{2\pi}
	\int d\nu'
	\mathcal K_{\alpha\beta}(\nu,\nu',\omega)
	Q_\beta(\nu',\omega)
	\label{eq:U_def}
\end{equation}
\begin{equation}
	V_{\alpha\beta}(\nu',\omega)
	=
	\frac{1}{2\pi}
	\int d\nu
	Q_\alpha(\nu,\omega)
	\mathcal K_{\alpha\beta}(\nu,\nu',\omega)
	\label{eq:V_def}
\end{equation}
The two-sided projected kernel is
\begin{equation}
	\widetilde{\mathcal K}_{\alpha\beta}(\omega)
	=
	\frac{1}{(2\pi)^2}
	\int d\nu\,d\nu'
	Q_\alpha(\nu,\omega)
	\mathcal K_{\alpha\beta}(\nu,\nu',\omega)
	Q_\beta(\nu',\omega)
	\label{eq:K_tilde_def}
\end{equation}
so that
\begin{align}
	\widetilde{\mathcal K}_{\alpha\beta}(\omega)
	&=
	\frac{1}{2\pi}
	\int d\nu\,
	Q_\alpha(\nu,\omega)
	U_{\alpha\beta}(\nu,\omega)
	\\
	&=
	\frac{1}{2\pi}
	\int d\nu'
	V_{\alpha\beta}(\nu',\omega)
	Q_\beta(\nu',\omega)
	\label{eq:K_tilde_one_sided}
\end{align}
At this stage, Eq.~\eqref{eq:K_tilde_def} is a definition of a projected kernel and is not yet an identity for the full BSE pole equation.

Define the integrated transition amplitude
\begin{equation}
	X_\alpha
	=
	\frac{1}{2\pi}
	\int d\nu\,
	\mathcal X_\alpha(\nu)
	\label{eq:X_def}
\end{equation}
and the projector onto the canonical $Q$ frequency shape
\begin{equation}
	(\mathcal P f)_\alpha(\nu)
	=
	Q_\alpha(\nu,\Omega)
	\frac{1}{2\pi}
	\int d\bar\nu\,
	f_\alpha(\bar\nu)
	\label{eq:P_def}
\end{equation}
with complementary projector
\begin{equation}
	\mathcal C
	=
	1-\mathcal P
	\label{eq:C_projector}
\end{equation}
The exact transition amplitude is decomposed as
\begin{equation}
	\mathcal X_\alpha(\nu)
	=
	Q_\alpha(\nu,\Omega)X_\alpha
	+
	\chi_\alpha(\nu)
	\label{eq:X_decomposition}
\end{equation}
with
\begin{equation}
	\frac{1}{2\pi}
	\int d\nu\,
	\chi_\alpha(\nu)
	=
	0
	\label{eq:chi_zero_average}
\end{equation}
The amplitude $\chi$ contains the frequency shapes discarded by the ordinary two-sided projection.

%%%%%%%%%%%%%%%%%%%%%%%%%
\section[Centered external vertices]{Centered external vertices and complementary frequency shapes}
\label{sec:centered_vertices}
%%%%%%%%%%%%%%%%%%%%%%%%%

Expand the kernel in the interaction
\begin{equation}
	\mathcal K
	=
	\mathcal K^{(1)}
	+
	\mathcal K^{(2)}
	+
	\mathcal K^{(3)}
	+
	\mathcal O(v^4)
	\label{eq:K_expansion}
\end{equation}
The first-order kernel is frequency independent.
It therefore maps the canonical $Q$ frequency shape into itself and annihilates every complementary shape with zero integral
\begin{align}
	\mathcal C\mathcal A^{(1)}\mathcal P
	&=
	0
	\\
	\mathcal P\mathcal A^{(1)}\mathcal C
	&=
	0
	\\
	\mathcal C\mathcal A^{(1)}\mathcal C
	&=
	0
	\label{eq:first_order_decoupling}
\end{align}
where
\begin{equation}
	\mathcal A^{(n)}
	=
	Q L_0\mathcal K^{(n)}
	\label{eq:A_def}
\end{equation}

At second order, define the centered one-sided vertices
\begin{equation}
	\delta U^{(2)}_{\alpha\beta}(\nu,\omega)
	=
	U^{(2)}_{\alpha\beta}(\nu,\omega)
	-
	\widetilde{\mathcal K}^{(2)}_{\alpha\beta}(\omega)
	\label{eq:delta_U}
\end{equation}
\begin{equation}
	\delta V^{(2)}_{\alpha\beta}(\nu,\omega)
	=
	V^{(2)}_{\alpha\beta}(\nu,\omega)
	-
	\widetilde{\mathcal K}^{(2)}_{\alpha\beta}(\omega)
	\label{eq:delta_V}
\end{equation}
which satisfy
\begin{equation}
	\frac{1}{2\pi}
	\int d\nu\,
	Q_\alpha(\nu,\omega)
	\delta U^{(2)}_{\alpha\beta}(\nu,\omega)
	=
	0
	\label{eq:delta_U_zero}
\end{equation}
\begin{equation}
	\frac{1}{2\pi}
	\int d\nu\,
	\delta V^{(2)}_{\alpha\beta}(\nu,\omega)
	Q_\beta(\nu,\omega)
	=
	0
	\label{eq:delta_V_zero}
\end{equation}
The leading feedback from the complementary frequency-shape space into the projected pole equation has the structure
\begin{equation}
	\begin{split}
	\Delta\widetilde{\mathcal K}^{(4)}_{\alpha\gamma}(\omega)
	={}&
	\sum_\beta
	(L_0)_\beta(\omega)
	\frac{1}{2\pi}
	\int d\nu\,
	Q_\beta(\nu,\omega)
	\\
	&\times
	\delta V^{(2)}_{\alpha\beta}(\nu,\omega)
	\delta U^{(2)}_{\beta\gamma}(\nu,\omega)
	\end{split}
	\label{eq:complement_feedback}
\end{equation}
The superscript in Eq.~\eqref{eq:complement_feedback} anticipates the order established below.

%%%%%%%%%%%%%%%%%%%%%%%%%
\section[Complementary-space feedback]{Audit of the complementary-space feedback}
\label{sec:feedback_audit}
%%%%%%%%%%%%%%%%%%%%%%%%%

The second-order one-sided vertices contain only $\hhp$ and $\pph$ poles.
For a forward intermediate pair $\beta=(jb,+)$, define
\begin{align}
	h^-_I(\nu)
	&=
	\frac{1}{\nu-E^-_I-\ii0^+}
	\\
	h^+_J(\nu)
	&=
	\frac{1}{\nu+\omega-E^+_J+\ii0^+}
	\label{eq:h_pm}
\end{align}
with
\begin{align}
	d^-_{\beta I}
	&=
	\omega-\eps_b+E^-_I
	\\
	d^+_{\beta J}
	&=
	\omega+\eps_j-E^+_J
	\label{eq:d_pm}
\end{align}
The causal prescriptions are understood in the following compact identities.
Using the $Q^+$ contour,
\begin{align}
	\langle h^-_I\rangle_\beta
	&=
	-\frac{1}{d^-_{\beta I}}
	\\
	\langle h^+_J\rangle_\beta
	&=
	+\frac{1}{d^+_{\beta J}}
	\label{eq:h_means}
\end{align}
where
\begin{equation}
	\langle f\rangle_\beta
	=
	\frac{1}{2\pi}
	\int d\nu\,
	Q^+_{jb}(\nu,\omega)f(\nu)
	\label{eq:average_def}
\end{equation}
Partial fractions give
\begin{equation}
	\left\langle
	\left(h^-_I-\langle h^-_I\rangle\right)
	\left(h^-_K-\langle h^-_K\rangle\right)
	\right\rangle_\beta
	=
	0
	\label{eq:cov_mm}
\end{equation}
\begin{equation}
	\left\langle
	\left(h^+_J-\langle h^+_J\rangle\right)
	\left(h^+_L-\langle h^+_L\rangle\right)
	\right\rangle_\beta
	=
	0
	\label{eq:cov_pp}
\end{equation}
Thus same-particle-number pole sectors do not couple through the complementary frequency-shape space.
For a mixed pair of poles,
\begin{equation}
	\begin{split}
	&\left\langle
	\left(h^-_I-\langle h^-_I\rangle\right)
	\left(h^+_J-\langle h^+_J\rangle\right)
	\right\rangle_\beta
	\\
	&\qquad=
	-
	\frac{
	\omega-(\eps_b-\eps_j)
	}{
	d^-_{\beta I}
	d^+_{\beta J}
	(\omega+E^-_I-E^+_J)
	}
	\end{split}
	\label{eq:cov_mp}
\end{equation}
The potentially resonant factor is therefore proportional to the inverse free-pair denominator.
Indeed,
\begin{equation}
	(L_0^+)_{jb}(\omega)
	\left\langle
	\delta h^-_I
	\delta h^+_J
	\right\rangle_\beta
	=
	-\ii
	\frac{1}
	{
	d^-_{\beta I}
	d^+_{\beta J}
	(\omega+E^-_I-E^+_J)
	}
	\label{eq:L0_cov_cancel}
\end{equation}
up to the infinitesimal causal prescriptions.
The apparent pole at $\omega=\eps_b-\eps_j$ cancels exactly.
The backward identity follows from Eq.~\eqref{eq:Q_reversal}.

Each residue of $\delta U^{(2)}$ and $\delta V^{(2)}$ is second order in the interaction.
Provided the $\hhp$, $\pph$, and mixed denominators in Eq.~\eqref{eq:L0_cov_cancel} remain non-intrusive in the perturbative sense, Eq.~\eqref{eq:complement_feedback} is therefore genuinely fourth order
\begin{equation}
	\Delta\widetilde{\mathcal K}_{\rm comp}
	=
	\mathcal O(v^4)
	\label{eq:feedback_order}
\end{equation}
This result is stronger than a naive order count because it explicitly removes the possible resonant enhancement of the intermediate $1h1p$ propagator.

%%%%%%%%%%%%%%%%%%%%%%%%%
\section[Pole-equation equivalence]{Equivalence of the pole equations through third order}
\label{sec:equivalence}
%%%%%%%%%%%%%%%%%%%%%%%%%

The preceding audit gives the central perturbative result.
Under the usual non-intruder assumptions, the full three-frequency pole equation and the two-sided projected BSE have identical poles through third order
\begin{equation}
	\text{three-frequency BSE}
	\overset{\mathcal O(v^3)}{\longleftrightarrow}
	\text{two-sided projected BSE}
	\label{eq:equivalence_statement}
\end{equation}
The projected equation is
\begin{equation}
	X
	=
	L_0(\omega)
	\left[
	\widetilde{\mathcal K}^{(1)}
	+
	\widetilde{\mathcal K}^{(2)}(\omega)
	+
	\widetilde{\mathcal K}^{(3)}(\omega)
	\right]
	X
	+
	\mathcal O(v^4)
	\label{eq:projected_pole_equation}
\end{equation}
The first omitted contribution is the centered $\hhp$--$\pph$ feedback of Eq.~\eqref{eq:complement_feedback}.
The statement is perturbative and assumes that no $N\pm1$ or mixed denominator becomes parametrically small with the interaction strength.

%%%%%%%%%%%%%%%%%%%%%%%%%
\section{Projected kernel through third order}
\label{sec:projected_third}
%%%%%%%%%%%%%%%%%%%%%%%%%

Separate the projected kernel into static and dynamic parts
\begin{equation}
	\widetilde{\mathcal K}^{[3]}(\omega)
	=
	\widetilde{\mathcal K}_{\rm stat}^{[3]}
	+
	\widetilde{\mathcal K}_{\rm dyn}^{[3]}(\omega)
	\label{eq:static_dynamic}
\end{equation}
All frequency-independent contributions can be absorbed into a physical Casida matrix $M_1^{[3]}$.
The dynamic terms inherit the three classes of the primitive multiresolvent expansion
\begin{equation}
	\widetilde{\mathcal K}_{\rm dyn}^{(2)}
	=
	C^{(1)}R^{(0)}D^{(1)}
	\label{eq:K_dyn_second}
\end{equation}
\begin{equation}
	\begin{split}
	\widetilde{\mathcal K}_{\rm dyn}^{(3)}
	={}&
	C^{(2)}R^{(0)}D^{(1)}
	+
	C^{(1)}R^{(0)}D^{(2)}
	\\
	&+
	C^{(1)}R^{(0)}M^{(1)}R^{(0)}D^{(1)}
	\end{split}
	\label{eq:K_dyn_third}
\end{equation}
The first two terms arise from second-order corrections to the two external couplings, while the third arises from a first-order interaction during the auxiliary propagation.
The matrices $C$ and $D$ introduced here are auxiliary couplings of the projected kernel and should not be confused with the primitive external vertices $U$ and $V$ before the second frequency integration.

%%%%%%%%%%%%%%%%%%%%%%%%%
\section{Second-order neutral-double pole space}
\label{sec:double_space}
%%%%%%%%%%%%%%%%%%%%%%%%%

The resonant $++$ block of the second-order GF2-type kernel can be written as
\begin{equation}
	\begin{split}
	[A_c^{(2)}(\omega)]_{ia,jb}
	={}&
	-
	\sum_{kc}
	\frac{
	\mel{jc}{}{ik}
	\mel{ka}{}{cb}
	}{
	\omega-(\eps_b-\eps_i+\eps_c-\eps_k)
	}
	\\
	&-
	\sum_{kc}
	\frac{
	\mel{jk}{}{ic}
	\mel{ca}{}{kb}
	}{
	\omega-(\eps_a-\eps_j+\eps_c-\eps_k)
	}
	\\
	&+
	\frac{1}{2}
	\sum_{kl}
	\frac{
	\mel{aj}{}{kl}
	\mel{lk}{}{bi}
	}{
	\omega-(\eps_a+\eps_b-\eps_k-\eps_l)
	}
	\\
	&+
	\frac{1}{2}
	\sum_{cd}
	\frac{
	\mel{aj}{}{cd}
	\mel{dc}{}{bi}
	}{
	\omega-(\eps_c+\eps_d-\eps_i-\eps_j)
	}
	\end{split}
	\label{eq:GF2_A}
\end{equation}
Every pole is a neutral-double energy
\begin{equation}
	\Delta_{kl}^{cd}
	=
	\eps_c+\eps_d-\eps_k-\eps_l
	\label{eq:Delta_double}
\end{equation}
For a unique double $I=(kl;cd)$ with $k<l$ and $c<d$, a convenient four-component attachment factorization is
\begin{equation}
	\mathbf C_{ia,I}
	=
	\begin{pmatrix}
	\delta_{ik}\mel{al}{}{cd}
	\\
	-\delta_{il}\mel{ak}{}{cd}
	\\
	-\delta_{ac}\mel{kl}{}{id}
	\\
	+\delta_{ad}\mel{kl}{}{ic}
	\end{pmatrix}
	\label{eq:C_attach}
\end{equation}
\begin{equation}
	\mathbf D_{I,jb}
	=
	\begin{pmatrix}
	\delta_{jk}\mel{dc}{}{lb}
	\\
	-\delta_{jl}\mel{dc}{}{kb}
	\\
	-\delta_{bc}\mel{jd}{}{kl}
	\\
	+\delta_{bd}\mel{jc}{}{kl}
	\end{pmatrix}
	\label{eq:D_attach}
\end{equation}
with attachment matrix
\begin{equation}
	S
	=
	\begin{pmatrix}
	0&1&1&1
	\\
	1&0&1&1
	\\
	1&1&0&1
	\\
	1&1&1&0
	\end{pmatrix}
	\label{eq:S_attach}
\end{equation}
The residue associated with the pole $I$ is
\begin{equation}
	\mathcal R^{(2)}_{ia,jb;I}
	=
	\mathbf C_{ia,I}^{\mathsf T}
	S
	\mathbf D_{I,jb}
	\label{eq:residue_attach}
\end{equation}
The eigenvalues of $S$ are $3,-1,-1,-1$, so the generic residue has rank four in attachment space.
One scalar auxiliary state per physical neutral double is therefore not sufficient in general.

A convenient generalized-resolvent realization uses
\begin{equation}
	\eta_c
	=
	\mathbf 1_{2h2p}\otimes S^{-1}
	\label{eq:eta_core}
\end{equation}
\begin{equation}
	M_c^{(0)}
	=
	\Delta_2\otimes S^{-1}
	\label{eq:M_core_zero}
\end{equation}
so that
\begin{equation}
	R_c^{(0)}(\omega)
	=
	[\omega\eta_c-M_c^{(0)}]^{-1}
	\label{eq:R_core_zero}
\end{equation}
contains
\begin{equation}
	[R_c^{(0)}(\omega)]_{I\lambda,J\mu}
	=
	\delta_{IJ}
	\frac{S_{\lambda\mu}}
	{\omega-\Delta_I+\ii0^+}
	\label{eq:R_core_elements}
\end{equation}
The backward sector follows by pair reversal and carries poles at $-\Delta_I$.
This four-copy realization is convenient but not unique; any exact rank factorization of the residues is equivalent.

%%%%%%%%%%%%%%%%%%%%%%%%%
\section[Neutral-double interaction]{First-order interaction inside the neutral-double space}
\label{sec:double_interaction}
%%%%%%%%%%%%%%%%%%%%%%%%%

The central third-order class is generated by a first-order interaction between neutral doubles.
For an antisymmetric amplitude $T_{ij}^{ab}$, define
\begin{equation}
	\begin{split}
	(H_2^{(1)}T)_{ij}^{ab}
	={}&
	\frac{1}{2}
	\sum_{cd}
	\mel{ab}{}{cd}
	T_{ij}^{cd}
	\\
	&+
	\frac{1}{2}
	\sum_{kl}
	\mel{kl}{}{ij}
	T_{kl}^{ab}
	\\
	&+
	\mathcal P(ij)\mathcal P(ab)
	\sum_{kc}
	\mel{kb}{}{cj}
	T_{ik}^{ac}
	\end{split}
	\label{eq:H2_action}
\end{equation}
The first line describes particle-particle scattering, the second hole-hole scattering, and the last line the four possible particle-hole scatterings.
The diagonal matrix element is
\begin{equation}
	\begin{split}
	(H_2^{(1)})_{ijab,ijab}
	={}&
	\mel{ab}{}{ab}
	+
	\mel{ij}{}{ij}
	\\
	&-
	\mel{ai}{}{ai}
	-
	\mel{aj}{}{aj}
	-
	\mel{bi}{}{bi}
	-
	\mel{bj}{}{bj}
	\end{split}
	\label{eq:H2_diagonal}
\end{equation}
which contains all six pairwise interactions among the four quasiparticles exactly once.
In the four-copy realization of Sec.~\ref{sec:double_space}, one convenient lift is
\begin{equation}
	M_c^{(1)}
	=
	H_2^{(1)}\otimes S^{-1}
	\label{eq:M_core_first}
\end{equation}
This choice reproduces the particle-particle, hole-hole, and particle-hole central third-order classes.
The auxiliary matrix itself is representation dependent; Eq.~\eqref{eq:M_core_first} is a convenient realization rather than a unique definition.

%%%%%%%%%%%%%%%%%%%%%%%%%
\section[External residues and spectators]{Third-order external residues and spectator labels}
\label{sec:spectator_space}
%%%%%%%%%%%%%%%%%%%%%%%%%

The external third-order classes generate single neutral-double poles with residues that are more general than the four attachment directions of Sec.~\ref{sec:double_space}.
A representative fully audited $+-$ contribution from the crossed $\eh$ primitive is
\begin{equation}
	\begin{split}
	\left[
	\ii\widetilde\Gamma^{\teh,(3)}_{\rm ext}
	\right]^{+-}_{ia,jb}(\omega)
	\supset{}&
	-
	\sum_{mekc}
	\frac{
	\mel{ac}{}{mk}
	\mel{ek}{}{jc}
	\mel{bm}{}{ie}
	}{
	[\omega-\Delta_{mk}^{ac}+\ii0^+]
	\Delta_{jm}^{ae}
	}
	\\
	&+
	\sum_{mekc}
	\frac{
	\mel{ak}{}{mc}
	\mel{ec}{}{jk}
	\mel{bm}{}{ie}
	}{
	[\omega+\Delta_{jk}^{ec}-\ii0^+]
	\Delta_{jm}^{ae}
	}
	\\
	&+
	\frac{1}{2}
	\sum_{mekl}
	\frac{
	\mel{ae}{}{kl}
	\mel{kl}{}{jm}
	\mel{bm}{}{ie}
	}{
	[\omega-\Delta_{kl}^{ae}+\ii0^+]
	\Delta_{jm}^{ae}
	}
	\\
	&-
	\frac{1}{2}
	\sum_{mecd}
	\frac{
	\mel{ae}{}{cd}
	\mel{cd}{}{jm}
	\mel{bm}{}{ie}
	}{
	[\omega+\Delta_{jm}^{cd}-\ii0^+]
	\Delta_{jm}^{ae}
	}
	\end{split}
	\label{eq:crossed_external_example}
\end{equation}
where
\begin{equation}
	\Delta_{ij}^{ab}
	=
	\eps_a+\eps_b-\eps_i-\eps_j
	\label{eq:Delta_general}
\end{equation}
The first pole in Eq.~\eqref{eq:crossed_external_example}, for example, may be factorized using a spectator-labelled auxiliary state
\begin{equation}
	\mathfrak I
	=
	(i\,|\,mk;ac)
	\label{eq:spectator_example}
\end{equation}
with
\begin{equation}
	C^{(1)}_{ia,\mathfrak I}
	=
	\mel{ac}{}{mk}
	\label{eq:spectator_C}
\end{equation}
\begin{equation}
	D^{(2)}_{\mathfrak I,jb}
	=
	-
	\sum_e
	\frac{
	\mel{ek}{}{jc}
	\mel{bm}{}{ie}
	}{
	\Delta_{jm}^{ae}
	}
	\label{eq:spectator_D}
\end{equation}
The spectator $i$ does not modify the pole energy $\Delta_{mk}^{ac}$; it retains an external-line degree of freedom required by the residue.

A corresponding audited contribution from the two-particle branch of the $\pp$ primitive is
\begin{equation}
	\begin{split}
	\left[
	\ii\widetilde\Gamma^{\pp,(3)}_{\ee,\rm ext}
	\right]^{+-}_{ia,jb}
	={}&
	-
	\sum_{\substack{c<d\\me}}
	\frac{
	\mel{ab}{}{cd}
	\mel{im}{}{ce}
	\mel{je}{}{dm}
	}{
	\Delta_{ij}^{cd}
	[\omega-\Delta_{im}^{ce}+\ii0^+]
	}
	\\
	&-
	\sum_{\substack{c<d\\me}}
	\frac{
	\mel{ab}{}{cd}
	\mel{ie}{}{dm}
	\mel{jm}{}{ce}
	}{
	\Delta_{ij}^{cd}
	[\omega+\Delta_{jm}^{ce}-\ii0^+]
	}
	\\
	&+
	\sum_{\substack{c<d\\me}}
	\frac{
	\mel{ab}{}{cd}
	\mel{ie}{}{cm}
	\mel{jm}{}{de}
	}{
	\Delta_{ij}^{cd}
	[\omega+\Delta_{jm}^{de}-\ii0^+]
	}
	\\
	&+
	\sum_{\substack{c<d\\me}}
	\frac{
	\mel{ab}{}{cd}
	\mel{im}{}{de}
	\mel{je}{}{cm}
	}{
	\Delta_{ij}^{cd}
	[\omega-\Delta_{im}^{de}+\ii0^+]
	}
	\end{split}
	\label{eq:pp_external_example}
\end{equation}
Equation~\eqref{eq:pp_external_example} uses a normalized unique central particle-pair basis $c<d$.
With unrestricted $c,d$ sums the corresponding explicit pair factor must be restored.
The two-hole pp branch and the opposite propagation orientations follow by the same pair-reversal operation used for the second-order external vertices.

Equations~\eqref{eq:crossed_external_example} and \eqref{eq:pp_external_example} show that the four-copy second-order attachment space is not generically closed under third-order external corrections.
A spectator-labelled construction gives an explicit polynomial residue factorization, while a minimal representation may instead be obtained by rank factorization of each pole residue.
For a pole $I$, the minimal number of auxiliary copies is
\begin{equation}
	r_I
	=
	\operatorname{rank}
	\mathcal R_I^{[3]}
	\label{eq:minimal_rank}
\end{equation}
The spectator construction is therefore a convenient exact realization rather than a claim of minimality.

%%%%%%%%%%%%%%%%%%%%%%%%%
\section[Frequency-independent linearization]{Frequency-independent linearization}
\label{sec:linearization}
%%%%%%%%%%%%%%%%%%%%%%%%%

For bookkeeping, split the auxiliary space into three sectors
\begin{equation}
	\mathcal Q
	=
	\mathcal Q_c
	\oplus
	\mathcal Q_L
	\oplus
	\mathcal Q_R
	\label{eq:Q_aux_split}
\end{equation}
The core sector $\mathcal Q_c$ contains the second-order neutral-double residues and the first-order internal interaction $M_c^{(1)}$.
The sector $\mathcal Q_L$ contains residue directions for which the left coupling first appears at second order, while $\mathcal Q_R$ contains the mirror class for which the right coupling first appears at second order.
Thus
\begin{align}
	C_c
	&=
	C_c^{(1)}
	\\
	D_c
	&=
	D_c^{(1)}
	\\
	C_L
	&=
	C_L^{(2)}
	\\
	D_L
	&=
	D_L^{(1)}
	\\
	C_R
	&=
	C_R^{(1)}
	\\
	D_R
	&=
	D_R^{(2)}
	\label{eq:coupling_orders}
\end{align}
All static terms through third order are collected in $M_1^{[3]}$.
The frequency-independent generalized eigenvalue problem is
\begin{equation}
	\begin{pmatrix}
	M_1^{[3]}
	& C_c^{(1)}
	& C_L^{(2)}
	& C_R^{(1)}
	\\
	D_c^{(1)}
	& M_c^{(0)}+M_c^{(1)}
	& 0
	& 0
	\\
	D_L^{(1)}
	& 0
	& M_L^{(0)}
	& 0
	\\
	D_R^{(2)}
	& 0
	& 0
	& M_R^{(0)}
	\end{pmatrix}
	\begin{pmatrix}
	X
	\\
	Y_c
	\\
	Y_L
	\\
	Y_R
	\end{pmatrix}
	=
	\omega
	\begin{pmatrix}
	\eta_1&0&0&0
	\\
	0&\eta_c&0&0
	\\
	0&0&\eta_L&0
	\\
	0&0&0&\eta_R
	\end{pmatrix}
	\begin{pmatrix}
	X
	\\
	Y_c
	\\
	Y_L
	\\
	Y_R
	\end{pmatrix}
	\label{eq:generalized_eigenproblem}
\end{equation}
Eliminating the auxiliary amplitudes gives
\begin{equation}
	\begin{split}
	M_{\rm eff}(\omega)
	={}&
	M_1^{[3]}
	+
	C_c^{(1)}R_c^{(0)}D_c^{(1)}
	\\
	&+
	C_L^{(2)}R_L^{(0)}D_L^{(1)}
	+
	C_R^{(1)}R_R^{(0)}D_R^{(2)}
	\\
	&+
	C_c^{(1)}R_c^{(0)}M_c^{(1)}R_c^{(0)}D_c^{(1)}
	+
	\mathcal O(v^4)
	\end{split}
	\label{eq:Schur_effective}
\end{equation}
which reproduces Eqs.~\eqref{eq:K_dyn_second} and \eqref{eq:K_dyn_third} term by term.
Therefore
\begin{equation}
	\begin{array}{c}
	\text{projected BSE}
	\\[-1mm]
	\Updownarrow\;\mathcal O(v^3)
	\\[-1mm]
	\text{generalized eigenvalue problem}
	\end{array}
	\label{eq:linearization_equivalence}
\end{equation}
Combining Eqs.~\eqref{eq:equivalence_statement} and \eqref{eq:linearization_equivalence} gives the desired third-order chain
\begin{equation}
	\begin{array}{c}
	\text{three-frequency BSE}
	\\[-1mm]
	\Updownarrow\;\mathcal O(v^3)
	\\[-1mm]
	\text{projected BSE}
	\\[-1mm]
	\Updownarrow\;\mathcal O(v^3)
	\\[-1mm]
	\text{frequency-independent enlarged eigenproblem}
	\end{array}
	\label{eq:full_chain}
\end{equation}
The matrices in Eq.~\eqref{eq:generalized_eigenproblem} are perturbatively constructed.
Diagonalizing the truncated matrix without further expansion generates partial higher-order resummations, exactly as in standard secular perturbative constructions.

%%%%%%%%%%%%%%%%%%%%%%%%%
\section{Construction recipe}
\label{sec:recipe}
%%%%%%%%%%%%%%%%%%%%%%%%%

The complete third-order construction can be organized as follows.

\begin{enumerate}
	\item Construct the direct $\eh$ irreducible vertex through third order from
	\begin{equation}
		\Gamma^\eh
		=
		\Lambda
		+
		\Phi^{\teh}
		+
		\Phi^\pp
	\end{equation}
	with $\Lambda=v+\mathcal O(v^4)$ and generate the transverse contribution by crossing

	\item Form the exact one-sided projections $U$ and $V$ using Eqs.~\eqref{eq:U_def} and \eqref{eq:V_def}

	\item Form the two-sided projected kernel $\widetilde{\mathcal K}$ using Eq.~\eqref{eq:K_tilde_def}

	\item If equivalence with the full three-frequency pole equation is required, center the second-order one-sided vertices according to Eqs.~\eqref{eq:delta_U} and \eqref{eq:delta_V} and verify the covariance identities of Sec.~\ref{sec:feedback_audit}

	\item Separate static and dynamic terms and absorb the static contribution through third order into $M_1^{[3]}$

	\item Collect each dynamic denominator by its neutral-double pole $\omega=\pm\Delta_{ij}^{ab}$

	\item Factor the second-order pole residues into the core auxiliary space $\mathcal Q_c$

	\item Construct the first-order neutral-double interaction $M_c^{(1)}$ from Eq.~\eqref{eq:H2_action}

	\item Factor the third-order single-pole residues generated by the two external classes into $\mathcal Q_L$ and $\mathcal Q_R$ using either spectator-labelled states or an exact rank factorization

	\item Assemble Eq.~\eqref{eq:generalized_eigenproblem} and expand its Schur complement through third order

	\item Verify pole positions, causal prescriptions, crossing, pair reversal, antisymmetry, unique-pair normalization, and the three perturbative classes in Eq.~\eqref{eq:K_dyn_third}
\end{enumerate}

The construction is representation independent: different factorizations of the same pole residues lead to algebraically equivalent enlarged eigenvalue problems.

%%%%%%%%%%%%%%%%%%%%%%%%%
\section{What changes at fourth order}
\label{sec:fourth_order}
%%%%%%%%%%%%%%%%%%%%%%%%%

The first contribution that is absent from the two-sided projected pole equation is the complementary-space feedback of Eq.~\eqref{eq:complement_feedback}.
After the cancellation in Eq.~\eqref{eq:L0_cov_cancel}, its characteristic denominator structure is
\begin{equation}
	\frac{1}
	{
	d^-_{\beta I}
	d^+_{\beta J}
	(\omega+E^-_I-E^+_J)
	}
	\label{eq:fourth_denominator}
\end{equation}
so the first genuinely new frequency-shape feedback mixes the $(N-1)$ $\hhp$ and $(N+1)$ $\pph$ sectors.
At the same perturbative order one also encounters the first correction to the fully irreducible vertex and higher corrections to the external and central auxiliary blocks.
Third order is therefore the highest order for which the full three-frequency pole problem reduces to the projected kernel without introducing the mixed complementary frequency-shape sector explicitly.

%%%%%%%%%%%%%%%%%%%%%%%%%
\section{Summary}
\label{sec:summary}
%%%%%%%%%%%%%%%%%%%%%%%%%

The exact three-frequency BSE is not generally reducible to one frequency by simple integration.
Through third order, however, centering the second-order one-sided vertices makes the same-sector $\hhp$ and $\pph$ covariances vanish, while the mixed covariance cancels the potentially resonant free $1h1p$ denominator.
Under non-intruder conditions, the full and projected pole equations are therefore perturbatively equivalent through third order.
The projected rational kernel contains neutral-double poles, a first-order interaction inside the neutral-double sector, and third-order external residues that may require spectator-labelled auxiliary directions.
Its Schur-complement linearization yields a frequency-independent generalized Casida/ADC-like eigenvalue problem through third order.

\end{document}
